\chapter{Lecture 7: Supersymmetry, Superfields, and Chirality in the One-Mode Toy Model}

These companion notes follow Leonard Susskind's lecture, with curation by LazyingArt LLC. The lecture proceeds by a deliberate retreat to the simplest possible supersymmetric system: one bosonic mode, one fermionic mode, both taken at rest, with all spatial dependence temporarily suppressed. From that toy model Susskind rebuilds the subject in stages. We first rewrite the bosonic mode as an ordinary harmonic oscillator, then discover why the fermionic Lagrangian must be first order in time, then review the supercharge algebra, and only after that pass to superspace, superfields, and the first recipe for constructing supersymmetric Lagrangians. The lecture is mathematically serious, but it is also exploratory; where a sign is corrected live, or where the blackboard calculation breaks off, the notes should preserve that rhythm.

\section{Bosons, Fermions, and the Zero-Momentum Toy System}

Susskind begins by saying, in effect, let us come back to fermions and bosons for a moment, and let us forget motion. The particles are either present or absent; momentum plays no role. In that setting the natural description is in terms of creation and annihilation operators. He uses the older \(+\) and \(-\) notation on the board, but in normalized form the Hamiltonian is
\begin{equation}
H = m\left(a^\dagger a + c^\dagger c\right),
\end{equation}
up to the usual ignored ground-state constant. Here \(a^\dagger,a\) create and annihilate the bosonic quantum, and \(c^\dagger,c\) do the same for the fermion.

This opening matters. We are not reviewing the algebra for its own sake. The point is to isolate the most elementary boson--fermion pairing before we add the machinery that will later guarantee that the pairing survives interactions. In this one-mode system the boson and fermion already have the same energy \(m\), and supersymmetry is the statement that there is an operation that exchanges them without changing that energy.

\section{The Bosonic Oscillator in Lagrangian Form}

Once the operator language is in place, Susskind changes viewpoint. Instead of talking about occupation numbers, he invents a field, or in this stripped-down context simply a coordinate \(\phi(t)\). The bosonic field theory with no space dependence is just a single harmonic oscillator.

\begin{figure}[h]
\centering
\includegraphics[width=0.62\textwidth]{lecture_07_figure_02.png}
\caption{Bosonic oscillator Lagrangian from the one-mode field setup. The board still shows the pre-correction \(T+V\) sign.}
\end{figure}

The board frame is worth keeping because it captures the exact transition: the operator-language Hamiltonian sits above, while the one-mode Lagrangian is written below. It also records a live correction. The board shows
\[
L=\frac{\dot\phi^{\,2}}{2}+\frac{m^2\phi^2}{2},
\]
but Susskind immediately stops himself: that is \(T+V\), the energy, whereas the Lagrangian must be \(T-V\). The corrected expression is therefore
\begin{equation}
L_b=\frac{1}{2}\dot\phi^{\,2}-\frac{1}{2}m^2\phi^2.
\end{equation}

Now the Euler--Lagrange equation is completely straightforward. We compute
\begin{align}
\frac{\partial L_b}{\partial \dot\phi} &= \dot\phi,
&
\frac{d}{dt}\left(\frac{\partial L_b}{\partial \dot\phi}\right) &= \ddot\phi,
\\
\frac{\partial L_b}{\partial \phi} &= -m^2\phi.
\end{align}
So
\begin{equation}
\ddot\phi + m^2\phi = 0.
\end{equation}
This is the ordinary oscillator equation, and the frequency is \(m\), not \(m^2\). That point is stated explicitly in the lecture. In this one-dimensional world the bosonic field is nothing but the familiar harmonic oscillator in disguise.

\section{Grassmann Variables and the Fermionic Lagrangian}

The fermionic part begins with the first real puzzle of the lecture. Susskind does not start by writing down the final fermionic Lagrangian. Instead he deliberately asks what happens if we try the simplest first-order expression. Take
\begin{equation}
L=\phi\,\dot\phi.
\end{equation}
If \(\phi\) is an ordinary bosonic variable, this does nothing:
\begin{equation}
L=\phi\,\dot\phi=\frac{1}{2}\frac{d}{dt}\phi^2.
\end{equation}
It is a total derivative, and the Euler--Lagrange equation is just the tautology \(\dot\phi=\dot\phi\). Every differentiable trajectory satisfies it.

Now comes the surprise. Suppose the same symbol denotes a fermionic variable, that is, a Grassmann variable. At this point the lecture renames it \(\psi\), precisely to avoid confusing bosons and fermions. We still have
\begin{equation}
L=\psi\,\dot\psi,
\qquad
\frac{\partial L}{\partial \dot\psi}=\psi.
\end{equation}
But when we differentiate with respect to \(\psi\), we must first move \(\psi\) past \(\dot\psi\), and that produces a minus sign:
\begin{equation}
\frac{\partial L}{\partial \psi}=-\,\dot\psi.
\end{equation}
So the equation of motion becomes
\begin{equation}
\dot\psi=-\,\dot\psi,
\end{equation}
hence
\begin{equation}
\dot\psi=0.
\end{equation}
This is no longer trivial. It is too strong for what we want, but it has real dynamical content.

\subsection{Question \& Answer}

\paragraph{Question.} Why can a first-order term be trivial for bosons but dynamical for fermions?

\paragraph{Answer.} For commuting variables, \(L=\phi\dot\phi\) is literally a total time derivative, so the Euler--Lagrange equation collapses to a tautology. For Grassmann variables, the same algebraic expression is sensitive to order. The variation with respect to \(\psi\) forces us to move \(\psi\) through \(\dot\psi\), and because both are anticommuting, that reordering contributes a minus sign. The lecture lingers over this point because it is the first moment when the Grassmann character of the variable does visible work in the equations of motion.

The story is not finished, because a single fermion is too naive. If we try to add a restoring term \(m\psi^2\), it vanishes identically: \(\psi^2=0\). This is exactly where Susskind repairs the model by passing to two real Grassmann variables \(\psi_1,\psi_2\), or equivalently to one complex fermion. In the lecture he first writes a mixed first-order Lagrangian for \(\psi_1\) and \(\psi_2\), with a mass term proportional to \(m\psi_1\psi_2\), and only afterwards repackages the result into complex notation.

The exact blackboard normalization wobbles, and the transcript is garbled for a short interval here, so the cleanest faithful reconstruction is to state the stabilized form of the answer:
\begin{equation}
L_f \sim \frac{i}{2}\left(\dot\psi^\ast \psi+\dot\psi\,\psi^\ast\right)+m\,\psi^\ast\psi.
\end{equation}
The important structural points are the ones Susskind emphasizes aloud. The kinetic term contains only one time derivative, and the restoring term contains one power of \(m\), not \(m^2\). The resulting equation of motion is first order:
\begin{equation}
i\,\dot\psi = m\,\psi.
\end{equation}
So the fermionic mode oscillates, just as the bosonic mode did, but now through a first-order equation rather than a second-order one. In this toy setting that is the one-dimensional shadow of the contrast between the Dirac equation and the Klein--Gordon equation.

Putting the bosonic and fermionic parts together, the lecture is aiming at a Lagrangian of schematic form
\begin{equation}
L \sim \frac{1}{2}\dot\phi^{\,2}-\frac{1}{2}m^2\phi^2
+\frac{i}{2}\left(\dot\psi^\ast \psi+\dot\psi\,\psi^\ast\right)
+m\,\psi^\ast\psi,
\end{equation}
with the explicit warning that the fermionic half-factors and ordering have been stabilized editorially rather than copied as exact board text.

\section{Supercharges and the Basic Supersymmetry Algebra}

With the toy bosonic and fermionic oscillators in hand, Susskind now returns to what supersymmetry means. The system has a symmetry that removes a fermion and inserts a boson, or vice versa, without changing the energy. The corresponding supercharges are
\begin{equation}
Q^\dagger=\sqrt{2m}\,a^\dagger c,
\qquad
Q=\sqrt{2m}\,a\,c^\dagger.
\end{equation}
The factor \(\sqrt{2m}\) is introduced purely by convention so that the algebra comes out in the standard form.

The key relation is the anticommutator:
\begin{equation}
\{Q,Q^\dagger\}=2m\left(a^\dagger a+c^\dagger c\right)=2H.
\end{equation}
The nilpotence is equally important:
\begin{equation}
Q^2=0,
\qquad
(Q^\dagger)^2=0,
\end{equation}
because every term in \(Q^2\) or \((Q^\dagger)^2\) contains the square of a fermionic creation or annihilation operator. And finally,
\begin{equation}
[H,Q]=[H,Q^\dagger]=0.
\end{equation}

This last relation is interpreted carefully in the lecture. It means, first, that supersymmetry is a genuine symmetry of the system. But it also means that the \(Q\)'s are conserved quantities. That is why they are called supercharges. Susskind spends a moment on this because it keeps the algebra from becoming empty symbolism: \([H,Q]=0\) means that the operation of trading a boson for a fermion is not merely imaginable, but compatible with the dynamics.

\section{From Supercharge Algebra to Superspace Geometry}

Up to this point the lecture has spoken algebraically. The next move is to reinterpret the same algebra geometrically. The guiding thought is familiar from ordinary quantum mechanics: generators of symmetry are differential operators. The Hamiltonian is the generator of time translations,
\begin{figure}[h]
\centering
\includegraphics[width=0.50\textwidth]{lecture_07_figure_03.png}
\caption{Hamiltonian as the generator of time translations.}
\end{figure}
\begin{equation}
H=i\,\frac{\partial}{\partial t},
\end{equation}
and momentum is the generator of spatial translations,
\begin{equation}
P=-\,i\,\frac{\partial}{\partial x}.
\end{equation}
Angular momentum is likewise built from coordinates and derivatives. So the natural question is: if \(Q\) and \(Q^\dagger\) are fermionic symmetry generators, might they be derivatives with respect to fermionic coordinates?

That is the entry into superspace. We enlarge the coordinate set from \(t\) to
\begin{equation}
(t,\theta,\bar\theta),
\end{equation}
where \(\theta\) is a Grassmann coordinate and \(\bar\theta\) its conjugate. The lecture is explicit that the sign conventions are unstable on the board. Susskind corrects himself, debates the signs with the room, and finally settles only on the fact that the two signs inside the pair of operators must match. We therefore choose one consistent convention:
\begin{equation}
q=\frac{\partial}{\partial\theta}+i\,\bar\theta\,\frac{\partial}{\partial t},
\qquad
\bar q=\frac{\partial}{\partial\bar\theta}+i\,\theta\,\frac{\partial}{\partial t}.
\end{equation}
Then
\begin{equation}
\{q,\bar q\}=2i\,\frac{\partial}{\partial t}=2H.
\end{equation}
This immediately explains why the time-derivative pieces are necessary. If \(q\) were only \(\partial_\theta\) and \(\bar q\) only \(\partial_{\bar\theta}\), the anticommutator would vanish. The extra \(i\bar\theta\partial_t\) and \(i\theta\partial_t\) pieces are forced on us by the algebra.

In this language a supersymmetry transformation is a tiny shift in the Grassmann coordinates together with a compensating shift in time:
\begin{equation}
\theta\to\theta+\xi,
\qquad
\bar\theta\to\bar\theta+\bar\xi,
\qquad
t\to t+i\left(\xi\bar\theta+\bar\xi\theta\right),
\end{equation}
again in a cleaned convention. The variation of a superfield \(\Phi\) can then be written as
\begin{equation}
\delta\Phi=\left(\xi q+\bar\xi\,\bar q\right)\Phi.
\end{equation}
That is the geometric heart of the lecture: supersymmetry is being re-read as motion in an enlarged space whose extra directions are Grassmann-valued.

\section{Superfields, Superactions, and the D-Term}

Once the geometric picture is in place, the lecture changes mood. The question is no longer what supersymmetry is, but how to build Lagrangians that keep it intact. For that we introduce a superfield.

In this one-dimensional superspace the most general bosonic superfield is finite, because \(\theta^2=\bar\theta^2=0\):
\begin{equation}
\Phi(t,\theta,\bar\theta)
=
\phi(t)+\bar\theta\,\psi(t)+\bar\psi(t)\,\theta+D(t)\,\bar\theta\theta.
\end{equation}
The termination of the expansion is not incidental; it is one of the practical beauties of Grassmann variables, and Susskind remarks on it more than once. The coefficient \(D(t)\) is another bosonic variable. At this stage it is simply the last coefficient in the power series.

Now we need an action. The superspace action is
\begin{equation}
S=\int dt\,d\theta\,d\bar\theta\,\Lambda,
\end{equation}
where \(\Lambda\) is a super-Lagrangian. This is the central constructive move of the lecture. One now builds \(\Lambda\) from superfields and their allowed derivatives, and then integrates over superspace just as one integrates an ordinary Lagrangian over ordinary space-time.

Berezin integration is simple but decisive:
\begin{equation}
\int d\theta\,d\bar\theta\,\Lambda=(\Lambda)_{\bar\theta\theta}.
\end{equation}
In words: integrating over \(\theta\) and \(\bar\theta\) picks out the coefficient of the highest Grassmann term. Susskind calls that surviving piece the \(D\)-term of the super-Lagrangian, though of course it need not be the auxiliary field \(D(t)\) itself. The ordinary Lagrangian is precisely this \(\bar\theta\theta\) coefficient.

That is why superspace is more than a bookkeeping device. It gives us a calculus for building candidate ordinary Lagrangians that are guaranteed, by construction, to fit inside a supersymmetric super-Lagrangian.

\section{Powers of Superfields, Derivative Subtleties, and Chirality}

Before confronting derivatives, the lecture takes up the easier operation: multiplication. The product of superfields is again a superfield, so the first nontrivial worked example is to square \(\Phi\).

\begin{figure}[h]
\centering
\includegraphics[width=0.82\textwidth]{lecture_07_figure_05.png}
\caption{Product of superfields expanded in components. At this point in the lecture the safest cleaned interpretation is the square \(\Phi^2\).}
\end{figure}

The board work is valuable here because Susskind actually multiplies the terms out, line by line, and then gathers the \(\bar\theta\theta\) coefficient. In cleaned form one finds
\begin{align}
\Phi^2
&=
\left(\phi+\bar\theta\psi+\bar\psi\theta+D\,\bar\theta\theta\right)^2
\\
&=
\phi^2
+2\phi\,\bar\theta\psi
+2\phi\,\bar\psi\theta
+\left(2\phi D-2\bar\psi\psi\right)\bar\theta\theta.
\end{align}
The ordering of the fermions has been normalized so that the bars remain on the left, exactly as Susskind says he prefers. Berezin integration then pulls out
\begin{equation}
\int d\theta\,d\bar\theta\,\Phi^2
=
2\phi D-2\bar\psi\psi.
\end{equation}
This is already a perfectly good supersymmetric contribution to the ordinary Lagrangian. More generally, the same logic applies to \(\Phi^3\), \(\Phi^4\), and indeed \(\Phi^n\): one computes only enough of the product to isolate the \(\bar\theta\theta\) term.

At this point the lecture turns sharply. The expression above is supersymmetric, but it is not yet interesting dynamically, because it contains no time derivatives. So the next question becomes: what kind of derivative of a superfield is still the right sort of superspace object?

\subsection{Question \& Answer}

\paragraph{Question.} When is a constraint on a field compatible with the symmetry?

\paragraph{Answer.} Susskind first answers this by an ordinary analogy. Suppose a field \(f(x,y)\) is constrained by
\begin{equation}
\frac{\partial f}{\partial y}=0.
\end{equation}
Then \(f\) depends only on \(x\). But if the theory is rotationally invariant, such a constraint is not stable under rotations, because the rotation generator \(L\) does not commute with \(\partial_y\); schematically,
\begin{equation}
[\partial_y,L]\sim \partial_x\neq 0.
\end{equation}
So the test is simple: a constraint is compatible with a symmetry only when the operator defining the constraint commutes with the symmetry generators. In the supersymmetric case the generators are fermionic, so the right condition is anticommutation rather than ordinary commutation.

That is exactly what motivates the covariant superspace derivatives. To avoid colliding with the auxiliary coefficient \(D(t)\), we write them as \(\mathcal D\) and \(\bar{\mathcal D}\):
\begin{equation}
\mathcal D=\frac{\partial}{\partial\theta}-i\,\bar\theta\,\frac{\partial}{\partial t},
\qquad
\bar{\mathcal D}=\frac{\partial}{\partial\bar\theta}-i\,\theta\,\frac{\partial}{\partial t}.
\end{equation}
These are built to anticommute with the supercharges:
\begin{equation}
\{\mathcal D,q\}=\{\mathcal D,\bar q\}=\{\bar{\mathcal D},q\}=\{\bar{\mathcal D},\bar q\}=0.
\end{equation}
Now one may impose the chiral constraint
\begin{equation}
\bar{\mathcal D}\Phi=0.
\end{equation}
Because \(\bar{\mathcal D}\) anticommutes with \(q\), chirality is preserved by a supersymmetry variation:
\begin{equation}
\bar{\mathcal D}(q\Phi)=-\,q(\bar{\mathcal D}\Phi)=0.
\end{equation}
That is the superspace version of the ordinary compatibility test.

The lecture then solves the constraint by introducing the shifted time coordinate
\begin{equation}
\tau=t+i\,\bar\theta\theta.
\end{equation}
A superfield that depends on \(\tau\) and \(\theta\), but not on \(\bar\theta\) independently, is automatically chiral:
\begin{equation}
\Phi(\tau,\theta)=\phi(\tau)+\bar\psi(\tau)\,\theta.
\end{equation}
Indeed,
\begin{equation}
\bar{\mathcal D}\Phi(\tau,\theta)
=
\left(\frac{\partial}{\partial\bar\theta}-i\,\theta\,\frac{\partial}{\partial t}\right)\Phi(\tau,\theta)
=
i\,\theta\,\frac{\partial\Phi}{\partial\tau}
-
i\,\theta\,\frac{\partial\Phi}{\partial t}
=0,
\end{equation}
since the dependence on \(t\) comes only through \(\tau\). Expanding back in \(t\),
\begin{align}
\phi(\tau)
&=
\phi\!\left(t+i\bar\theta\theta\right)
=
\phi(t)+i\,\bar\theta\theta\,\dot\phi(t),
\end{align}
with no higher term because \((\bar\theta\theta)^2=0\). Thus the chiral superfield takes the explicit form
\begin{equation}
\Phi(t,\theta,\bar\theta)
=
\phi(t)+\bar\psi(t)\,\theta+i\,\bar\theta\theta\,\dot\phi(t).
\end{equation}

Here the lecture reaches its most fragile point. Susskind now turns to the product of a chiral superfield with its complex conjugate, because that is the natural route to the kinetic super-Lagrangian. But at the board he begins to lose confidence in the ordering and the signs. He starts the conjugate expression, hesitates, and finally says that he is too tired to trust the derivation. That stopping point matters. We should not smooth it away. The lecture has earned the right to stop exactly where the speaker stops.

\section{Summary}

The chapter has unfolded in the same order as the lecture. We began with bosons and fermions at rest, described only by occupation numbers, and then rewrote the bosonic mode as an ordinary harmonic oscillator. We then met the first genuinely supersymmetric surprise: a first-order Lagrangian that is trivial for a commuting variable becomes dynamical for a Grassmann one. Repairing that naive fermionic model led us back to the matched bosonic and fermionic oscillators and to the supercharges that exchange them.

From there the lecture shifted perspective. The algebra of \(Q\) and \(Q^\dagger\) was reinterpreted geometrically by enlarging time to superspace, with Grassmann coordinates \(\theta\) and \(\bar\theta\). Once that enlargement was in place, the problem became constructive: define superfields, define a superaction, use Berezin integration to extract the ordinary Lagrangian, and then build candidate invariant terms from powers of the superfield. The last step was subtler. Ordinary time derivatives were not enough; one needed covariant derivatives and then a constraint, \(\bar{\mathcal D}\Phi=0\), that was compatible with the symmetry. This produced the chiral superfield and brought the lecture to the edge of the kinetic construction.

The final punchline is the one Susskind insists on, even after the algebra becomes exhausting. A supersymmetric Lagrangian is not merely a convenient way of packaging fields. It enforces a matching between bosons and fermions strong enough to produce exact cancellation statements. That is why the machinery matters. At the same time, the lecture ends honestly: the clean \(\bar\Phi\Phi\) computation is deferred, and the later breaking of supersymmetry is left for another discussion. The notes should preserve both facts.